\section{Adapter tuning for NLP}

\begin{SCfigure*}
\begin{tabular}{cc}
 \includegraphics[width=0.45\linewidth]{figures/Adapter_insertion.pdf}&
 \includegraphics[width=0.45\linewidth]{figures/Adapter_arch.pdf}
\end{tabular}
 \caption{
 Architecture of the adapter module and its integration with the Transformer.
 \textbf{Left:} We add the adapter module twice to each Transformer layer:
 after the projection following multi-headed attention and after the two feed-forward layers.
 \textbf{Right:} The adapter consists of a bottleneck which contains few parameters relative to the attention and feedforward layers in the original model.
 The adapter also contains a skip-connection.
 During adapter tuning, the green layers are trained on the downstream data, this includes the adapter, the layer normalization parameters,
 and the final classification layer (not shown in the figure).
 \label{fig:adapters_transformer}}
\end{SCfigure*}

We present a strategy for tuning a large text model on several downstream tasks.
Our strategy has three key properties:
(i) it attains good performance,
(ii) it permits training on tasks sequentially, that is, it does not require simultaneous access to all datasets,
and (iii) it adds only a small number of additional parameters per task.
These properties are especially useful in the context of cloud services,
where many models need to be trained on a series of downstream tasks, so a high degree of sharing is desirable.

To achieve these properties, we propose a new bottleneck adapter module.
Tuning with adapter modules involves adding a small number of new parameters to a model, which are trained on the downstream task~\citep{rebuffi2017}.
When performing vanilla fine-tuning of deep networks, a modification is made to the top layer of the network.
This is required because the label spaces and losses for the upstream and downstream tasks differ.
Adapter modules perform more general architectural modifications to re-purpose a pre-trained network for a downstream task.
In particular, the adapter tuning strategy involves injecting new layers into the original network.
The weights of the original network are untouched, whilst the new adapter layers are initialized at random.
In standard fine-tuning, the new top-layer and the original weights are co-trained.
In contrast, in adapter-tuning, the parameters of the original network are frozen and therefore may be shared by many tasks.

Adapter modules have two main features: a small number of parameters, and a near-identity initialization.
The adapter modules need to be small compared to the layers of the original network.
This means that the total model size grows relatively slowly when more tasks are added.
A near-identity initialization is required for stable training of the adapted model;
we investigate this empirically in Section~\ref{sec:discussion}.
By initializing the adapters to a near-identity function, original network is unaffected when training starts.
During training, the adapters may then be activated to change the distribution of activations throughout the network.
The adapter modules may also be ignored if not required;
in Section~\ref{sec:discussion} we observe that some adapters have more influence on the network than others.
We also observe that if the initialization deviates too far from the identity function, the model may fail to train.

\subsection{Instantiation for Transformer Networks\label{sec:bottleneckadapter}}

We instantiate adapter-based tuning for text Transformers.
These models attain state-of-the-art performance in many NLP tasks,
including translation, extractive QA, and text classification problems~\citep{vaswani2017,radford2018improving,devlin2018bert}.
We consider the standard Transformer architecture, as proposed in~\citet{vaswani2017}.

Adapter modules present many architectural choices.
We provide a simple design that attains good performance.
We experimented with a number of more complex designs, see Section~\ref{sec:discussion},
but we found the following strategy performed as well as any other that we tested, across many datasets.

Figure~\ref{fig:adapters_transformer} shows our adapter architecture, and its application it to the Transformer.
Each layer of the Transformer contains two primary sub-layers: an attention layer and a feedforward layer.
Both layers are followed immediately by a projection that maps the features size back to the size of layer's input.
A skip-connection is applied across each of the sub-layers.
The output of each sub-layer is fed into layer normalization.
We insert two serial adapters after each of these sub-layers.
The adapter is always applied directly to the output of the sub-layer, after the projection back to the input size,
but before adding the skip connection back.
The output of the adapter is then passed directly into the following layer normalization.

To limit the number of parameters, we propose a bottleneck architecture.
The adapters first project the original $d$-dimensional features into a smaller dimension, $m$, apply a nonlinearity, then project back to $d$ dimensions.
The total number of parameters added per layer, including biases, is $2md+d+m$.
By setting $m\ll d$, we limit the number of parameters added per task;
in practice, we use around $0.5-8\%$ of the parameters of the original model.
The bottleneck dimension, $m$, provides a simple means to trade-off performance with parameter efficiency.
The adapter module itself has a skip-connection internally.
With the skip-connection, if the parameters of the projection layers are initialized to near-zero,
the module is initialized to an approximate identity function.

Alongside the layers in the adapter module, we also train new layer normalization parameters per task.
This technique, similar to conditional batch normalization~\citep{de2017modulating},
FiLM~\citep{perez2018}, and self-modulation~\citep{chen2019}, also yields parameter-efficient adaptation of a network; with only $2d$ parameters per layer.
However, training the layer normalization parameters alone is insufficient for good performance,
see Section~\ref{sec:param_efficiency}.
